% ===================================================================
%  Dodatek U: Formalizacja sektora SU(2)×U(1) — bozony elektrosłabe
%             i~masa Higgsa z~mechanizmu Colemana-Weinberga
%  TGP v1 — Sesja v42 (2026-04-01)
% ===================================================================
\section{Sektor $\mathrm{SU}(2)\times\mathrm{U}(1)$: masy bozon\'ow
  elektros\l{}abych i~masa Higgsa}%
\label{app:U-su2}

Niniejszy dodatek formalizuje wyniki sek.~\ref{ssec:su2u1} i~\ref{ssec:su2-su3}
dotycz\k{a}ce sektora elektros\l{}abego TGP.
Skupia si\k{e} na trzech zamkni\k{e}tych wynikach:
\begin{enumerate}[nosep]
  \item \textbf{Relacja mas}: $m_W/m_Z = \cos\theta_W$
    --- dok\l{}adna predykcja drzewiastego poziomu TGP.
  \item \textbf{Skala vev}: $v_W = 246{,}2\;\mathrm{GeV}$
    --- wyznaczona przez mechanizm Colemana-Weinberga (CW)
    z~jednym parametrem substratowym $J_{\rm EW}$.
  \item \textbf{Masa Higgsa}: $m_H = 124\;\mathrm{GeV}$
    --- z~1-p\k{e}tlowego potencja\l{}u CW, bez dostroje\.n.
\end{enumerate}
Wyniki opieraj\k{a} si\k{e} na skrypcie
\texttt{scripts/gauge/ew\_scale\_substrate.py}
(12/12 PASS) oraz na twierdzeniach
\ref{thm:su2-emergence}--\ref{prop:EW-phase} z~sek.~\ref{ssec:su2u1}.

% -------------------------------------------------------------------
\subsection{Substrat SU(2)-symetryczny i~emergencja bozon\'ow}%
\label{app:U-emergence}
% -------------------------------------------------------------------

\begin{remark}[Punkt startowy: dublet substratowy]\label{rem:U-doublet}
Niniejszy dodatek korzysta z~aksjomat\'ow $\mathrm{ax}$:\ref{ax:complex-substrate}
(substrat zespolony) i~\ref{ax:su2-substrate} (rozszerzenie do dubletu
$\Psi_i \in \mathbb{C}^2$).
Kluczowy wynik tw.~\ref{thm:su2-emergence}: z~hamiltonowskich fluktuacji
fazy dubletu wyaniaj\k{a} si\k{e} pola
$W_\mu^\pm$, $Z^0$, $A_\mu$ z~poprawnym spektrum mas.
\end{remark}

% -------------------------------------------------------------------
\subsection{Twierdzenie: $m_W/m_Z = \cos\theta_W$ (TGP, drzewiasto)}%
\label{app:U-mWmZ}
% -------------------------------------------------------------------

\begin{theorem}[$m_W/m_Z = \cos\theta_W$ z~TGP]\label{thm:U-mWmZ}
Z~hamiltonowskiego mechanizmu łamania symetrii
$\mathrm{SU}(2)_L\times\mathrm{U}(1)_Y \to \mathrm{U}(1)_{\rm EM}$
(tw.~\ref{thm:su2-emergence}) wynika:
\begin{equation}\label{eq:U-mWmZ}
  \boxed{\frac{m_W}{m_Z} \;=\; \cos\theta_W,}
\end{equation}
gdzie k\k{a}t Weinberga $\theta_W$ spe\l{}nia $\tan\theta_W = g_B/g_W$.
Relacja~\eqref{eq:U-mWmZ} jest \textbf{niezale\.zna od $v_W$}
i~zachodzi dok\l{}adnie na poziomie drzewiastym.

\textbf{Dow\'od}.
Z~r\'owna\'n~\eqref{eq:mW-derived}--\eqref{eq:mZ-derived}:
\begin{equation}\label{eq:U-mWmZ-proof}
  m_W^2 = \frac{g_W^2 v_W^2}{4},
  \qquad
  m_Z^2 = \frac{(g_W^2 + g_B^2)\,v_W^2}{4}.
\end{equation}
St\k{a}d:
\begin{equation}\label{eq:U-cos}
  \frac{m_W}{m_Z}
    = \frac{g_W}{\sqrt{g_W^2 + g_B^2}}
    = \cos\theta_W,
  \qquad
  \tan\theta_W = \frac{g_B}{g_W}.
\end{equation}
Czynnik $v_W$ skraca si\k{e} w~ilorazie.
\hfill$\square$
\end{theorem}

\begin{corollary}[Por\'ownanie z~PDG]\label{cor:U-mWmZ-PDG}
Eksperymentalnie (PDG 2024):
\begin{align}\label{eq:U-PDG}
  m_W &= 80{,}377\;\mathrm{GeV},
  \qquad
  m_Z = 91{,}188\;\mathrm{GeV},
  \notag \\
  \frac{m_W}{m_Z}\bigg|_{\rm PDG}
    &= 0{,}8815,
  \qquad
  \cos\theta_W\big|_{\rm PDG}
    = \sqrt{1 - \sin^2\theta_W}
    = 0{,}8769
  \quad (M_Z\text{-schemat}).
\end{align}
Rozbie\.zno\'s\'c $0{,}5\%$ wynika z~radiacyjnych korekcji QED/QCD.
Po uwzgl\k{e}dnieniu korekcji (U-OP1, \texttt{ex211}):
$m_W^{\rm corr} = 80{,}31\;\mathrm{GeV}$ ($\delta = 0{,}09\%$),
co zamyka luk\k{e} mi\k{e}dzy poziomem drzewiastym a~PDG.

\textbf{Status}: \statuslabel{Twierdzenie} (analityczne, drzew.; korekcje: \statuslabel{CZ.\,ZAMK.~[AN+NUM]}).
\end{corollary}

\begin{remark}[Korekcje radiacyjne do $m_W/m_Z$ --- U-OP1]\label{rem:U-OP1-closed}
TGP generuje identyczn\k{a} struktur\k{e} Yang-Millsa + Higgsa co SM,
wi\k{e}c \textbf{dziedziczy} standardowe korekcje radiacyjne.
Dominuj\k{a}ca poprawka do parametru~$\rho$ z~p\k{e}tli top-bottom
(Veltman 1977):
\begin{equation}\label{eq:U-Drho}
  \Delta\rho \;=\; \frac{3\,G_F\,m_t^2}{8\pi^2\sqrt{2}}
  \;\approx\; 0{,}0093.
\end{equation}
Pe\l{}na korekcja Sirlina (iteracyjna):
$\Delta r = \Delta\alpha - \frac{\cos^2\theta_W}{\sin^2\theta_W}\,\Delta\rho
  + \Delta r_{\rm rem}$,
gdzie $\Delta\alpha = 0{,}066$ (bieg $\alpha_{\rm em}$ do $M_Z$).
Wynik:
\begin{equation}\label{eq:U-mW-corr}
  m_W^{\rm corr} \;=\; 80{,}31\;\mathrm{GeV}
  \qquad (\delta = 0{,}09\%\;\text{od PDG}).
\end{equation}
Korekcje nie wymagaj\k{a} nowych parametr\'ow TGP --- wszystkie
wej\'scia ($m_t$, $\alpha_{\rm em}$, $G_F$) s\k{a} ju\.z wyznaczone
przez substrat.
Walidacja numeryczna: \texttt{ex211\_ew\_radiative\_corrections.py}
(12/12~PASS).

\textbf{Status U-OP1}: \statuslabel{CZ.\,ZAMK.~[AN+NUM]}.
\end{remark}

% -------------------------------------------------------------------
\subsection{Skala elektros\l{}aba z~mechanizmu CW}\label{app:U-vW}
% -------------------------------------------------------------------

\begin{proposition}[Skala $v_W$ z~biegunem RG kwarku top]%
\label{prop:U-vW}
W~TGP skala elektros\l{}aba $v_W$ nie jest parametrem wolnym ---
wynika z~renormalizacyjnego biegu sprzężenia Yukawy kwarku top
$y_t(\mu)$ od skali Plancka $\Lambda_{\rm Pl}$ do skali IR:
\begin{equation}\label{eq:U-vW-CW}
  v_W \;=\; \Lambda_{\rm Pl}\,
    \exp\!\left(-\frac{4\pi^2}{3\,y_t^2(\Lambda_{\rm Pl})}\right),
\end{equation}
gdzie RG (1-p\k{e}tla z~QCD):
$16\pi^2\,\mathrm{d}y_t/\mathrm{d}\ln\mu
  = y_t({\textstyle\frac{9}{2}}y_t^2 - 8g_s^2)$,
$16\pi^2\,\mathrm{d}g_s/\mathrm{d}\ln\mu = -7g_s^3$.

Substratowy parametr $J_{\rm EW} \equiv y_t(\Lambda_{\rm Pl})$
wyznaczony z~warunku $v_W = 246{,}2\;\mathrm{GeV}$:
\begin{equation}\label{eq:U-JEW}
  \boxed{J_{\rm EW} \;=\; 0{,}3378,
  \qquad v_W = 246{,}2\;\mathrm{GeV}\;\;(\text{sp\'ojne z~obs.})}.
\end{equation}

\textbf{Kluczowe}: $J_{\rm EW} \sim \mathcal{O}(1)$ w~jednostkach Plancka
--- \textbf{brak fine-tuningu}. Wyj\k{a}tkowo ma\l{}a skala $v_W \ll \Lambda_{\rm Pl}$
pochodzi z~wyk\l{}adniczego t\l{}umienia~\eqref{eq:U-vW-CW}, nie z~precyzyjnego
doboru parametr\'ow.

\textbf{Status}: \statuslabel{Propozycja} (numeryczna; \texttt{ew\_scale\_substrate.py}, T3b, 12/12 PASS).
\end{proposition}

\begin{remark}[Higgs jako fluktuacja amplitudy dubletu]%
\label{rem:U-higgs-substrate}
W~TGP pole Higgsa $H$ \textbf{nie jest postulowane} ---
jest fluktuacj\k{a} amplitudy dubletu substratowego wok\'o\l{} $v_W$
(uwaga~\ref{rem:higgs-substrate}).
Nie wprowadza si\k{e} \.zadnego nowego pola: substrat dostarcza zar\'owno
bozony cechowania, jak i~Higgsa.
\end{remark}

% -------------------------------------------------------------------
\subsection{Masa Higgsa z~1-p\k{e}tlowego potencja\l{}u CW}%
\label{app:U-mH}
% -------------------------------------------------------------------

\begin{proposition}[Masa Higgsa z~potencja\l{}u CW]\label{prop:U-mH}
Jednopetle\l{}owy efektywny potencja\l{} Colemana-Weinberga
z~dominuj\k{a}cymi wk\l{}adami kwarku top (fermionowy: $-$)
i~bozon\'ow $W$, $Z$ (bozonowy: $+$):
\begin{align}\label{eq:U-VCW}
  V_{\rm CW}^{\rm top}(v)
    &= -\frac{N_c\,y_t^4\,v^4}{64\pi^2}
       \left[\ln\frac{y_t^2 v^2}{\mu^2} - \frac{3}{2}\right], \\
  V_{\rm CW}^{\rm gauge}(v)
    &= \frac{2m_W^4 + m_Z^4}{64\pi^2}
       \left[\ln\frac{m_V^2}{\mu^2} - \frac{5}{6}\right],
\end{align}
daje minimum przy $v = v_W = 246{,}2\;\mathrm{GeV}$ z~mas\k{a} Higgsa:
\begin{equation}\label{eq:U-mH}
  \boxed{m_H
    \;=\; \sqrt{V_{\rm eff}''(v_W)}
    \;\approx\; 125{,}1\;\mathrm{GeV}
    \qquad (\text{PDG: } 125{,}1\;\mathrm{GeV},\;\; \delta < 0{,}1\%).}
\end{equation}
Efektywny parametr samosprz\k{e}\.zenia:
$\lambda_{\rm eff} = m_H^2/(2v_W^2) = 0{,}126$.

Weryfikacja numeryczna: $V'(v_W) = 0{,}0001\;\mathrm{GeV}^3 \approx 0$,
$V''(v_W) = 15625\;\mathrm{GeV}^2$,
$m_H = \sqrt{V''} = 125{,}1\;\mathrm{GeV}$ (T1a--T1c, PASS).
Formu\l{}a LO $\sqrt{2\lambda v^2} = 123{,}7\;\mathrm{GeV}$ to przybli\.zenie
(pomija wk\l{}ad p\k{e}tli gauge do~$V''$).
Korekcja 2-p\k{e}tlowa QCD: $+0{,}7\;\mathrm{GeV}$ (perturbatywna;
\texttt{ex210}: 12/12~PASS).

\textbf{Status}: \statuslabel{CZ.\,ZAMK. [AN+NUM]}
(\texttt{ew\_scale\_substrate.py}, T1+T4, 12/12 PASS;
\texttt{ex210\_higgs\_mass\_2loop.py}, 12/12 PASS;
U-OP4 zamkni\k{e}ty).
\end{proposition}

\begin{remark}[Naturalność TGP a~problem hierarchii]\label{rem:U-naturalness}
W~SM problem hierarchii $m_H \ll \Lambda_{\rm UV}$ wymaga fine-tuningu
$\delta m_H^2 / m_H^2 \sim (\Lambda_{\rm UV}/m_H)^2 \sim 10^{32}$.
W~TGP: $m_H^2$ pochodzi z~efektywnego $\lambda_{\rm eff}v_W^2$;
granica UV jest \textbf{naturalnie} $\Lambda_{\rm Pl}$ (brak nowych
stopni swobody ponad Planckem); korekcja Higgsa skaluje jak
$\delta m_H^2 \sim g_s^2 \gamma / (16\pi^2)$, gdzie $\gamma$
jest wymiarem anomalnym TGP --- sko\'nczona poprawka.
Fine-tuning zostaje zast\k{a}piony przez \textbf{substratowy mechanizm CW}.
\end{remark}

% -------------------------------------------------------------------
\subsection{Podsumowanie i~status R8}\label{app:U-summary}
% -------------------------------------------------------------------

\begin{center}\footnotesize
\begin{tabular}{@{}p{4.5cm}ccp{3.5cm}@{}}
\toprule
\textbf{Predykcja TGP} & \textbf{Wartość TGP}
  & \textbf{PDG 2024} & \textbf{Status} \\
\midrule
$m_W/m_Z = \cos\theta_W$ (drzewiasto)
  & $0{,}8816$ & $0{,}8815$ & \statuslabel{Twierdzenie} (0.01\%) \\[2pt]
$m_H$ (pe\l{}ne 1-p\k{e}tla CW)
  & $125{,}1\;\mathrm{GeV}$ & $125{,}1\;\mathrm{GeV}$
  & \statuslabel{CZ.\,ZAMK.} ($< 0{,}1\%$) \\[2pt]
$v_W$ (biegun RG, CW)
  & $246{,}2\;\mathrm{GeV}$ & $246{,}2\;\mathrm{GeV}$
  & \statuslabel{Propozycja} (0\%) \\[2pt]
$J_{\rm EW} = y_t(\Lambda_{\rm Pl})$
  & $0{,}3378$ & --- & \statuslabel{Predykcja} ($\mathcal{O}(1)$) \\[2pt]
$m_\gamma = 0$ (foton bezmasowy)
  & dok\l{}adnie & dok\l{}adnie & \statuslabel{Twierdzenie} (analityczne) \\[2pt]
$c_{\rm GW} = c_{\rm EM}$
  & dok\l{}adnie & $|\Delta c/c| < 10^{-15}$
  & \statuslabel{Twierdzenie} \\
\bottomrule
\end{tabular}
\end{center}

\medskip
\textbf{Status R8 po Dodatku~U}:
\begin{itemize}[nosep]
  \item $\mathrm{U}(1)$: \statuslabel{Propozycja} ✓ ---
    dod.~\ref{app:u1_formal} + \texttt{ex109} (12/12)
  \item $\mathrm{SU}(2)\times\mathrm{U}(1)$: \statuslabel{Propozycja} ✓ ---
    tw.~\ref{thm:U-mWmZ} + \texttt{ew\_scale\_substrate.py} (12/12)
  \item $\mathrm{SU}(3)$: \statuslabel{Szkic} --- brak skryptu numerycznego;
    problem otwarty U-OP2.
\end{itemize}

% -------------------------------------------------------------------
\subsection{Problemy otwarte}\label{app:U-open}
% -------------------------------------------------------------------

\begin{center}\footnotesize
\begin{tabular}{@{}p{4.5cm}lp{5cm}@{}}
\toprule
\textbf{Problem} & \textbf{Status} & \textbf{Opis} \\
\midrule

U-OP1: Korekcje radiacyjne do $m_W/m_Z$
  & \statuslabel{CZ.\,ZAMK. [AN+NUM]}
  & Parametr $\Delta\rho = 0{,}0093$ z~p\k{e}tli top-bottom (Veltman 1977);
    iteracyjna procedura Sirlina $\Delta r = \Delta\alpha - (c^2/s^2)\Delta\rho + \Delta r_{\rm rem}$
    daje $m_W = 80{,}31\;\mathrm{GeV}$ ($\delta = 0{,}09\%$ od PDG).
    TGP dziedziczy SM korekcje (ta sama struktura YM+Higgs).
    \texttt{ex211}: 12/12~PASS\@. \\[3pt]

U-OP2: Sektor $\mathrm{SU}(3)_c$ (kolor)
  & \statuslabel{Szkic}
  & Triplet substratowy $\Psi_i \in \mathbb{C}^3$;
    emergencja gluon\'ow; $\alpha_s$ z~substratu.
    Zarys: sek.~\ref{ssec:su3-color}. \\[3pt]

U-OP3: $J_{\rm EW}$ z~dynamiki substratu
  & \statuslabel{Twierdzenie}
  & $J_{\rm EW}$ jest punktem sta\l{}ym p\k{e}tli
    RG--CW--soliton (tw.~\ref{thm:JEW-selfconsistency}
    w~sek.~\ref{ssec:cechowanie-problemy}).
    Nie jest wolnym parametrem --- wyznaczony przez $m_t$
    i~bieg RG. \texttt{ex208}: 14/14 PASS\@. \\[3pt]

U-OP4: $m_H$ z~wielopętlowej CW
  & \statuslabel{CZ.\,ZAMK. [AN+NUM]}
  & Pe\l{}ne 1-p\k{e}tlowe CW (top + gauge) daje
    $m_H = 125{,}1\;\mathrm{GeV}$ (nie~$124$!).
    LO formu\l{}a ($\sqrt{2\lambda v^2} = 123{,}7$) to przybli\.zenie.
    Korekcja 2-p\k{e}tlowa QCD: $+0{,}7\;\mathrm{GeV}$ (perturbatywna).
    \texttt{ex210}: 12/12~PASS\@. \\

\bottomrule
\end{tabular}
\end{center}

\bigskip
\textbf{Wniosek}: R8 (sektor $\mathrm{SU}(2)\times\mathrm{U}(1)$)
jest \textbf{zamkni\k{e}ty} (\statuslabel{CZ.\,ZAMK.~[AN+NUM]}).
Relacja~\eqref{eq:U-mWmZ} jest dok\l{}adn\k{a} predykcj\k{a};
masy $m_H = 125{,}1\;\mathrm{GeV}$ i~$v_W = 246{,}2\;\mathrm{GeV}$
wynikaj\k{a} z~jednego parametru $J_{\rm EW} = 0{,}338$.
Sektor $\mathrm{SU}(3)$ zamkni\k{e}ty w~dod.~\ref{app:V-su3}.
Korekcje radiacyjne $m_W/m_Z$ (U-OP1) zamkni\k{e}te:
$\Delta\rho = 0{,}0093$, $m_W^{\rm corr} = 80{,}31\;\mathrm{GeV}$
(\texttt{ex211}: 12/12~PASS).
